%!TEX root = ../main.tex

% Pendiente: actualizar los resultados de modelado (algoritmos seleccionados,
% desempeño y variables más influyentes) al cerrar las Etapas 4 y 5.

\begin{abstract}[it]
    La calidad del servicio de energía eléctrica en Colombia se evalúa mediante los indicadores regulatorios SAIDI (Índice de Duración Promedio de Interrupción del Sistema) y SAIFI (Índice de Frecuencia Promedio de Interrupción del Sistema), definidos por la CREG y vigilados por la SSPD. En la empresa de distribución objeto de estudio, ubicada en Norte de Santander, la información operativa se encuentra dispersa en plataformas heterogéneas (BRAE, OMS SPARD, OMS SP7 y GIS), con esquemas de codificación distintos y sin trazabilidad homogénea entre eventos y activos de red, lo que limita la construcción de modelos predictivos confiables e interpretables. El objetivo de este trabajo es desarrollar una metodología, basada en el estándar CRISP-DM, para pronosticar la frecuencia y la duración de las interrupciones eléctricas por zona geográfica ---variables proxy asociadas a SAIFI y SAIDI, respectivamente---, e identificar los factores técnicos, operativos y climáticos que explican su comportamiento. Se integraron seis fuentes corporativas y la fuente climática externa NASA POWER mediante consultas SQL y procesamiento en Python y R, conformando un panel de $N = 87{.}007{.}074$ observaciones con grano activo eléctrico $\times$ día ($A = 47{.}649$ activos $\times$ $T = 1{.}826$ días, 2021--2025) y una tabla agregada de $90{.}128$ observaciones con grano cuadrante $\times$ semana sobre una malla territorial de $800$ celdas. El análisis exploratorio aplicó diagnóstico de valores faltantes por variable y por fila, detección de valores extremos mediante la regla de Tukey (factores $1{.}5$ y $3{.}0$) y percentiles P1--P99, caracterización de distribuciones asimétricas y cero-infladas, análisis de series semanales con ACF y PACF, correlación de Spearman bivariada y parcial, coeficiente de Cramér, factor de inflación de la varianza (VIF), análisis de componentes principales (PCA), agrupamiento jerárquico de variables, un modelo lineal generalizado de Poisson exploratorio y segmentación territorial mediante \textit{k}-means. Los resultados de esta fase muestran que la trazabilidad evento--activo pasa de cerca del $31$~\% en 2021--2023 al $79{.}8$~\% en 2024 y al $95{.}0$~\% en 2025 tras la migración a OMS SP7, por lo que el universo de modelado se restringe a 2024--2025 ($36{.}120$ observaciones) con partición temporal de entrenamiento en 2024 y prueba en 2025; que la señal climática contemporánea sobre la frecuencia de interrupciones es débil en el grano cuadrante~$\times$~semana y está encabezada por el viento ($\hat{\rho} = 0{.}14$; $p < 0{.}001$) por delante de la precipitación acumulada semanal ($\hat{\rho} = 0{.}07$, que asciende a $\hat{\rho} = 0{.}11$ al controlar temperatura y viento), lo que traslada el peso predictivo hacia la estructura temporal y territorial de la propia serie de eventos; que existe fuerte redundancia entre temperaturas y vientos (VIF de hasta $43{.}92$, reducible a un máximo de $7{.}21$ con un subconjunto de cuatro covariables) y entre las variables operativas de severidad; y que la actividad de interrupciones presenta marcada heterogeneidad espacial, con el $14{.}2$~\% de los cuadrantes concentrando la mayor parte de los eventos. Con base en estos hallazgos se definen \texttt{n\_eventos} y \texttt{duracion\_total} como variables respuesta y un conjunto reducido de predictoras climáticas y operativas. En la etapa siguiente se compararán modelos estadísticos y de aprendizaje automático aptos para datos de conteo con exceso de ceros y autocorrelación temporal, evaluados mediante MAE, RMSE y MAPE junto con análisis de importancia de variables. Los productos del trabajo son un repositorio analítico trazable, informes reproducibles del análisis exploratorio y lineamientos para anticipar la degradación de la continuidad del servicio, priorizar mantenimiento preventivo en zonas críticas y apoyar el cumplimiento regulatorio en beneficio de más de $600{.}000$ usuarios de la región.

    \vspace{0.8em}
    \noindent\textbf{Palabras clave:} continuidad del servicio eléctrico; SAIDI; SAIFI; CRISP-DM; integración de datos; análisis exploratorio de datos; modelos predictivos; variables climáticas.
\end{abstract}
